\begin{conclusions}
\label{ch:conclusion}

本文面向通信资源受限的多无人机协同观测，以任务可靠完成为总体目标，研究了能量约束下的无人机路由与功率控制、消息长度预算约束下的语义信息选择，以及确认不及时条件下的发送与重传控制。三项研究分别改善转发能力、观测判断质量与实际接收结果，主要工作和结论如下。

（1）针对能量约束下吞吐量与网络寿命难以兼顾的问题，提出 LEADR 分布式路由与功率控制方法。用李雅普诺夫漂移加惩罚将长期能量要求写入时隙目标，根据链路速率与能量代价求解功率，再结合剩余能量、地理进度和运动预测选择下一跳。理论部分给出了功率子问题的求解条件，以及虚拟队列与累计预算的关系。节点数量、通信范围和飞行速度实验表明，LEADR 在保持较高转发效率与较低丢包的同时维持了任务周期内的通信能力，其收益来自服务量与能量使用的协调，而非单独最小化能量差异。

（2）针对有限通信量下重复或低质量观测占用发送机会的问题，提出 EviComRL 语义消息选择方法。使用参数固定的视觉语言模型生成候选，策略结合中心已有判断、链路代价、来源可靠性和选择历史决定发送内容。训练时排除超出预算和重复选择的动作，经策略预训练与 PPO 学习消息顺序和停止时机。基准配置下，分类准确率为58.17\%，各类别 F1 的平均值为50.58\%，每任务平均通信量为49.70 token。实验表明，观测覆盖影响可获得的证据，粒度和预算影响分类质量及通信量；链路恶化时，即使预测类别变化较小，有效信念质量也可能下降。可靠性特征的重复实验表明，可观测信念和节点共识能够辅助识别有用证据。

（3）针对数据到达与确认返回不一致的问题，提出 C-BFC 交付控制方法。对同一任务中预测类别相同的非空消息集合统一收益估计，根据累计确认和发送记录更新接收概率。控制器比较两步内新发、重传、等待和停止的预期收益，并按累计能耗调节后续任务支出。在17种链路与资源配置下，平均按时正确完成率为49.84\%，单位正确任务能耗为3.615 mJ。相较于 SD-SAC 和 TP-DQN，能耗分别降低16.59\%和14.32\%，平均正确完成率略有提高。消融结果显示，接收信念减少了将确认缺失当作接收失败的错误，等待有助于调整后续选择，一致性修正主要减少低收益通信。
\end{conclusions}
